% CLAUDE REPORT::WCDC-0001::2026-11-12T06:51-08:00
% !TEX program = xelatex
% Compiles with XeLaTeX or pdfLaTeX (ASCII source).
\documentclass[11pt]{article}
\usepackage[margin=1in]{geometry}
\usepackage{amsmath,amssymb,amsthm}
\usepackage{xurl}
\newtheorem{theorem}{Theorem}
\newtheorem{proposition}[theorem]{Proposition}
\newtheorem{lemma}[theorem]{Lemma}
\newtheorem{corollary}[theorem]{Corollary}
\theoremstyle{remark}
\newtheorem*{remark}{Remark}
\newcommand{\SAT}{\mathrm{SAT}}
\newcommand{\UNSAT}{\mathrm{UNSAT}}
\newcommand{\poly}{\operatorname{poly}}
\newcommand{\cost}{\operatorname{cost}}
\newcommand{\Wbar}{\overline W}
\newcommand{\status}[1]{\textbf{[#1]}}

\title{CLAUDE WCDC-0001\\[2pt]
\large Worst-case derivation complexity: an exact classification of the\\
bridge $\SAT\in\mathsf P\Rightarrow \Wbar_R(N)\le\poly(N)$}
\author{CLAUDE (scheduled independent run; red-team orientation)}
\date{12 November 2026, 06:51 PST}

\begin{document}
\maketitle

\section*{Status}
\textbf{P versus NP remains OPEN.} This report proves elementary but exact
facts about the shared worst-case objective and its proposed bridge.
Every theorem below is \status{PROVED/EXACT} or \status{KNOWN THEOREM} as
labelled; no load-bearing arrow of the project is closed. The main
finding is that the bridge is never "weaker" than the goal in the way the
directive hopes: for every fixed calculus $R$ it is \emph{either} logically
equivalent to $\mathsf P\neq\mathsf{NP}$ (when a lower bound for $R$ is known)
\emph{or} equivalent to the open statement
``$\mathsf P=\mathsf{NP}$ implies $R$ is p-optimal'' (Theorem~2).

\section{Objects and the cost hypothesis}
Fix a Cook--Reckhow \emph{refutation} system $R$ for $\UNSAT$ 3-CNFs: a
polynomial-time map from strings onto encodings of unsatisfiable 3-CNFs
(the multisign DAG hypergraph inference system of the directive is one
such $R$ once its rules are fixed and checkable in polynomial time). Set
\[
 w_R(F)=\min_{\pi:F\vdash_R\bot}\cost_R(\pi),\qquad
 \Wbar_R(N)=\max_{\substack{F\in\UNSAT\\|\langle F\rangle|\le N}} w_R(F).
\]
\textbf{(H\textsubscript{cost})} There is a fixed $c$ with
$|\pi|^{1/c}-c\le \cost_R(\pi)\le c\,(|\pi|+1)^c$, so cost and encoded length are
polynomially equivalent. (This is the comparison the corpus requires before
using Cook--Reckhow; the project's unit-node cost satisfies the upper side, and
the lower side holds whenever each node has polynomial-size encoding.)
Under (H\textsubscript{cost}), ``$\Wbar_R\le\poly$'' is the same as
``$R$ is polynomially bounded'' (p-bounded).

$R$ is \emph{optimal} if every refutation system $S$ satisfies
$w_R(F)\le\poly(w_S(F)+|F|)$, and \emph{p-optimal} if a polynomial-time
translation of $S$-refutations into $R$-refutations also exists.
Write
\[
 B_R:\quad \SAT\in\mathsf P\ \Longrightarrow\ \Wbar_R(N)\le\poly(N),
 \qquad
 L_R:\quad \Wbar_R(N)\ \text{is superpolynomial}.
\]
$B_R\wedge L_R$ gives $\mathsf P\neq\mathsf{NP}$ immediately (modus tollens).

\section{Results}

\begin{lemma}[free search is exact under the hypothesis]\status{PROVED/EXACT}
Assume $\mathsf P=\mathsf{NP}$ and that $R$ is p-bounded with bound $p$. Then a
minimum-length $R$-refutation of a given $F$ can be produced in time
$\poly(|F|)$.
\end{lemma}
\begin{proof}
``There is $\pi$ with $|\pi|\le k$, $|\pi|\le p(|F|)$ and $R(\pi)=F$'' is an
$\mathsf{NP}$ predicate in $(F,1^k)$. Under $\mathsf P=\mathsf{NP}$ decide it for
$k=0,1,\dots,p(|F|)$, then recover the bits of $\pi$ one at a time by the
usual search-to-decision self-reduction. Every call is polynomial in $|F|$.
\end{proof}
\begin{remark}
This vindicates the directive's clarification from the other side: not
charging strategy discovery costs nothing in the bridge direction. It also
explains why hardness of \emph{finding} short proofs (e.g.\ NP-hardness of
automating Resolution, Atserias--M\"uller 2019/2020) has no bearing on
$\Wbar_R$: those theorems concern search, which $\Wbar_R$ ignores by design.
\status{KNOWN THEOREM} for the cited result; the irrelevance is
\status{PROVED/EXACT}.
\end{remark}

\begin{theorem}[bridge = p-optimality conditional]\status{PROVED/EXACT}
Under (H\textsubscript{cost}), for every fixed refutation system $R$,
\[
 B_R\iff\bigl(\mathsf P=\mathsf{NP}\Rightarrow R\text{ is p-optimal}\bigr)
 \iff\bigl(\mathsf P=\mathsf{NP}\Rightarrow R\text{ is optimal}\bigr).
\]
\end{theorem}
\begin{proof}
Assume $\mathsf P=\mathsf{NP}$. Then $\mathsf{NP}=\mathsf{coNP}$, so a p-bounded
refutation system $Q$ exists (Cook--Reckhow 1979, \status{KNOWN THEOREM}).
(i) If $R$ is p-bounded, it is p-optimal: given an $S$-refutation $w$, compute
$F=S(w)$ (so $|F|\le\poly(|w|)$) and apply Lemma~1. (ii) p-optimal implies
optimal trivially. (iii) If $R$ is optimal, it size-simulates $Q$, so
$w_R(F)\le\poly(w_Q(F)+|F|)\le\poly(|F|)$, i.e.\ $R$ is p-bounded.
Hence under $\mathsf P=\mathsf{NP}$ the three properties coincide, and by
(H\textsubscript{cost}) p-boundedness is $\Wbar_R\le\poly$.
\end{proof}

\begin{corollary}[no slack anywhere]\status{PROVED/EXACT}
(a) If $L_R$ is already known, then $B_R\iff\mathsf P\neq\mathsf{NP}$.
This applies to Resolution, Res($k$), Cutting Planes, Polynomial Calculus,
bounded-depth Frege, and every system with a published superpolynomial
lower bound \status{KNOWN THEOREM for those lower bounds}.
(b) If $R\le_p R'$ ($R'$ p-simulates $R$) then $B_R\Rightarrow B_{R'}$ while
$L_{R'}\Rightarrow L_R$. So moving $R$ up makes the bridge more plausible
and the lower bound harder, and vice versa. The only systems where neither
half is already settled or equivalent to the goal are those whose lower
bound is open (Frege, Extended Frege and stronger), and there the bridge is
the classical open question whether $\mathsf P=\mathsf{NP}$ would force such a
system to be p-optimal.
\end{corollary}
\begin{proof}
(a) If $L_R$ holds, $B_R$ says ``$\SAT\in\mathsf P\Rightarrow$ false''. (b)
Simulation transports upper bounds up and lower bounds down.
\end{proof}

\begin{proposition}[adaptive strategy switching is a min]\status{PROVED/EXACT}
Let $R_1,\dots,R_k$ be refutation systems with $k$ fixed or $k\le\poly(N)$, and
let $R_\oplus$ accept $(i,\pi)$ iff $R_i$ accepts $\pi$. Then
$w_{R_\oplus}(F)=\min_i w_{R_i}(F)+O(\log k)$ and
\[
 \Wbar_{R_\oplus}(N)=\max_F\min_i w_{R_i}(F)+O(\log k)\ \le\ \min_i\Wbar_{R_i}(N)+O(\log k),
\]
where the minimum is taken \emph{inside} the maximum.
\end{proposition}
\begin{proof}
The verifier for $R_\oplus$ reads the $O(\log k)$-bit tag and runs $R_i$, so
$R_\oplus$ is a refutation system and the identity for $w$ is immediate. Taking
the maximum over $F$ gives the equality, and $\max_F\min_i\le\min_i\max_F$ gives the
inequality. Consequently a family that is hard for one $R_i$ does not lower-bound
$R_\oplus$ unless the same formulas are hard for every $R_i$. Example: mod-2 Tseitin
formulas on expanders are exponential for Resolution (Urquhart 1987;
Ben-Sasson--Wigderson 2001) but have polynomial $\mathrm{PC}_{\mathbb F_2}$
refutations (Gaussian elimination), so they give no lower bound for
$\mathrm{Res}\oplus\mathrm{PC}_{\mathbb F_2}$.
\end{proof}
\begin{remark}[consequence for the hard-family role]
A candidate hard family survives ``alternative reasoning methods'' only if the
\emph{same} formulas are hard for every system in the portfolio. Random 3-CNF
near the threshold is the standard test: exponential for Resolution
(Chv\'atal--Szemer\'edi 1988) and for Polynomial Calculus (Ben-Sasson--Impagliazzo;
Alekhnovich--Razborov) \status{KNOWN THEOREM}; Cutting Planes lower bounds are
known for random $k$-CNF with $k=\Theta(\log n)$ (Fleming--Pankratov--Pitassi--Robere
2017) \status{KNOWN THEOREM}, and the 3-CNF Cutting Planes case should be checked
against current literature before it is used. Tseitin is \status{FALSIFIED/RETIRED}
as a portfolio-hard family (easy over $\mathbb F_2$; quasipolynomial in Cutting
Planes, Dadush--Tiwari 2020).
\end{remark}

\begin{proposition}[what is known for $R=$ Extended Frege]\status{KNOWN THEOREM / CONDITIONAL}
(1) Cook 1975: if a polynomial-time $A$ decides $\SAT$ and $\mathrm{PV}$ proves
$\forall F,a\,(\mathrm{Sat}(F,a)\to A(F)=1)$, then EF is p-bounded. So $B_{EF}$
holds for every solver whose soundness is feasibly provable.
(2) Pich--Santhanam (arXiv:2312.08163, ECCC TR23-199): if $S^1_2$ proves that a
polynomial-time function generates anticheckers for SAT, or witnesses an error
of every polynomial-size circuit failing to solve SAT, then ``EF not p-bounded''
implies $\mathsf P\neq\mathsf{NP}$. These hypotheses are \status{OPEN}.
(3) The corpus reflection results (LSPV-1/336, /361; Codex 2026-11-11 transfer
audit) reduce $B_E$ for the project fragment $E$ to short $E$-refutations of the
solver's own reflection formulas; by Theorem~2 these are exactly the
p-optimality obligations, so the routes agree.
\end{proposition}

\section{Hostile controls}
\begin{itemize}
\item \textbf{Solver-induced system.} If $A$ is a polynomial SAT algorithm,
$S_A$ (proof of $F$ = $F$ itself, verifier runs $A$) is p-bounded. Any
lower bound $L_R$ therefore fails for $S_A$; Theorem~2 says $B_R$ is exactly the
demand that $R$ simulate $S_A$. The directive's lower bound must be proved for
a system that does not contain such clocks, and the bridge must then put them
back. This is the whole difficulty, not a technicality.
\item \textbf{Relativization.} Theorem~2 relativizes. So does Cook--Reckhow.
Neither is a barrier violation, and neither proves anything about P vs NP
alone.
\item \textbf{World ledger.} $B_R$ for a lower-bounded $R$ is a row of type
``equivalent of $\mathsf P\neq\mathsf{NP}$'' (a reformulation, not progress).
$B_{EF}$ is a ``world separator''-type open implication.
\end{itemize}

\section{Corpus reviewed}
Root \texttt{main.tex} (surveyed; LSPV ownership regions not edited because
this run is not an LSPV role); \texttt{AGENTS.md}; the newest commits through
LSPV-4/562 (06:48 PST) and the coordination handoffs of 11--12 November;
Codex \texttt{2026-11-11\_1203\_SAT\_POLYNOMIAL\_ASSUMPTION\_AUDIT.tex} and
\texttt{2026-11-11\_1331\_TRANSFER\_RECOVERY\_AND\_RATE\_AMPLIFICATION.tex}
(reflection, size optimality, Pich--Santhanam section); repository search for
LSPV-1 reflection/optimality reports 336, 337, 353, 361; the CLAUDE
world-ledger and earlier CLAUDE reports.

\section{Duplication / contradiction audit}
Corollary~3(a) duplicates the Codex 11-11 remark that, once the lower bound is
proved, the bridge is equivalent to the separation; it is restated for
completeness. Theorem~2 (bridge $\iff$ conditional p-optimality, with the three
notions collapsing under $\mathsf P=\mathsf{NP}$) and Lemma~1 as a justification of
free search were not found in the corpus as stated; the Codex audit notes that
size optimality plus a lower bound gives $\mathsf{NP}\ne\mathsf{coNP}$, which is the
converse direction. No contradiction with any verified result was found.

\section{What changed relative to root main}
Nothing in root \texttt{main.tex} is changed. Suggested for LSPV-5 synthesis:
record that the bridge target is exactly ``$\mathsf P=\mathsf{NP}\Rightarrow R$
p-optimal'', and that the hard-family role must test one family against the whole
portfolio (Proposition~4).

\section{Adopted / rejected / retired components}
Adopted: Cook--Reckhow, Cook 1975, Pich--Santhanam as cited. Rejected: any reading
of the bridge as strictly easier than $\mathsf P\neq\mathsf{NP}$ for a lower-bounded
system. Retired: Tseitin as a portfolio-hard family.

\section{Stale or overstated claims found}
None in the verified corpus. The scheduled directive's phrase ``weakest sufficient
bridge'' is accurate only relative to the pointwise simulation targets; by
Corollary~3 it carries the full difficulty.

\section{Literature checked this run}
Cook--Reckhow 1979; Cook 1975 (PV); Kraj\'{\i}\v{c}ek--Pudl\'ak 1989 (optimal
systems); Atserias--M\"uller, \emph{Automating Resolution is NP-hard} (FOCS 2019,
JACM 2020); Pich--Santhanam, \url{https://arxiv.org/abs/2312.08163},
\url{https://eccc.weizmann.ac.il/report/2023/199}. A search for 2025--2026 results
settling whether $\mathsf P=\mathsf{NP}$ implies EF is p-bounded found none.

\section{Next target}
Find the weakest feasibility hypothesis (in the sense of Pich--Santhanam) that a
\emph{fixed} combined inference system $R$ of the directive could satisfy, and check
whether adding its witnessing axioms keeps every lower bound already proved for $R$.

\end{document}
% END_OF_VERIFIED_REPORT::CLAUDE::WCDC-0001
